% GENERATED FILE. Edit canonical lessons, then run npm run generate:books.
% canonical-source-hash: fnv1a64:b6104ce0fb027d67
% canonical-lessons: ZH-C251-pian, ZH-C251-zuowen, ZH-W251-zuo-do, ZH-C251-xiezuo, ZH-W251-xie-write, ZH-C251-shengzi, ZH-C251-xiugai

\chapter{Essays and New Characters}
\label{ch:zh-essays-and-new-characters}
\hlchaptermodality{251}

\begin{chapteropening}
\textbf{By the end of this chapter:} \emph{I can count essays and stories, talk about a school essay, writing and new characters, and put mistakes right.}

Complete the last lesson of chapter 251: I can count essays and stories, talk about a school essay, writing and new characters, and put mistakes right.
\end{chapteropening}

\section[piān]{piān — measure word for essays and stories}

\label{lesson:ZH-C251-pian}



\begin{quote}
Before the new one: say the Mandarin for a story, then the Mandarin for to play chess.
\end{quote}

\subsection*{You'll want to know: piān}
\emph{piān} — \textquotedblleft{}measure word for essays and stories\textquotedblright{}.

\begin{quote}
\emph{zhè piān xiǎoshuō hěn cháng} — This story is a long one.
\end{quote}

Its characters come later; for now it is a word for the ear.

\begin{cousinweb}
\emph{piān} is inherited from Middle Chinese \emph{phjien} (Guangyun, in Baxter's transcription), glossed as a section of a book.
\end{cousinweb}

\begin{grammarlens}[title={What you've built}]
One more word for saying what I do.
\end{grammarlens}

\subsection*{Your turn}
\begin{itemize}
\raggedright
  \item \emph{Say it:} \emph{piān}
  \item \emph{Say it:} \emph{piān}, once more
  \item \emph{Say it:} \emph{zhè piān xiǎoshuō hěn cháng} — This story is a long one.
  \item \emph{Recall:} say the Mandarin for a story, then the Mandarin for to play chess, then say \emph{piān} again
\end{itemize}

\begin{tcolorbox}[breakable,colback=teal!4,colframe=teal!35!black,title={Before you move on}]
What does piān mean? (\textquotedblleft{}Measure word for essays and stories\textquotedblright{}.) Say it once more.
\end{tcolorbox}
\section[\texorpdfstring{\zh{作文} (zuòwén)}{作文 (zuòwén)}]{\zh{作文} (zuòwén) — essay; composition (at school)}

\label{lesson:ZH-C251-zuowen}



\begin{quote}
Before the new one: say the Mandarin for to play chess, then the Mandarin for measure word for essays.
\end{quote}

\subsection*{You'll want to know: \zh{作文}}
\textbf{\zh{作文}} — \emph{zuòwén} — \textquotedblleft{}essay; composition (at school)\textquotedblright{}.

\begin{quote}
\emph{wǒ zài xiě zuòwén} — I'm writing an essay.
\end{quote}

The new character, in pieces: \zh{作} — seven strokes: the person piece \zh{亻}; then a short stroke falling to the left, a bar, a long vertical, and two short bars to the right of it. The next lesson puts it in your hand.

\begin{cousinweb}
\textbf{\zh{作文}} (\emph{zuòwén}) is built from \zh{作} 'to make' and \zh{文} 'writing', both inherited from Middle Chinese, \emph{tsak} and \emph{mjun} (Guangyun, in Baxter's transcription), which glosses the first 'to do' and 'to rise', and the second 'a piece of writing'.
\end{cousinweb}

\begin{grammarlens}[title={What you've built}]
One more word for saying what I do.
\end{grammarlens}

\subsection*{Your turn}
\begin{itemize}
\raggedright
  \item \emph{Say it:} \emph{zuòwén}
  \item \emph{Say it:} \emph{zuòwén}, once more
  \item \emph{Say it:} \emph{wǒ zài xiě zuòwén} — I'm writing an essay.
  \item \emph{Recall:} say the Mandarin for to play chess, then the Mandarin for measure word for essays, then say \emph{zuòwén} again
\end{itemize}

\begin{tcolorbox}[breakable,colback=teal!4,colframe=teal!35!black,title={Before you move on}]
What does \zh{作文} mean? (\textquotedblleft{}Essay; composition (at school)\textquotedblright{}.) Say it once more.
\end{tcolorbox}
\section[\texorpdfstring{\zh{作} (zuò)}{作 (zuò)}]{\zh{作} — the person piece, then a post with two short bars}

\label{lesson:ZH-W251-zuo-do}



\begin{quote}
The word first, then the character inside it.

\begin{itemize}
\raggedright
  \item \emph{Recall:} say \emph{zuòwén} — \textbf{R1}, one lesson back, before it decays
\end{itemize}

\begin{quote}
\zh{做} \zh{你}
\end{quote}

\emph{Read it:} each one aloud, with nothing said first — \emph{zuò}, \emph{nǐ}
\end{quote}

\subsection*{Script — the person piece, then a post with two short bars}
\hlblockfigure{\detokenize{figures/ZH-W251-zuo-do-filmstrip.pdf}}{How \zh{作} is written, stroke by stroke}

\begin{quote}
\zh{作}
\end{quote}

\textbf{Two pieces, left before right.}

\textbf{The left piece, \zh{亻}, two strokes,} as in \zh{你} and \zh{做}: a stroke falling to the left, then a vertical down to the foot.

\textbf{The right piece, five strokes.} A short stroke falling to the left at the top. A bar to the right from its middle. A long vertical down from the bar to the foot. Then two short bars against the right side of the vertical, the upper one first.

2 + 5 = seven.

\subsection*{Your turn}
\begin{itemize}
\raggedright
  \item \emph{Write it:} \zh{作}, all seven, with the model in front of you
  \item \emph{Cover:} write it again from memory, and count the pieces, not the strokes
\end{itemize}

\begin{tcolorbox}[breakable,colback=teal!4,colframe=teal!35!black,title={Before you move on}]
How many strokes in \zh{作}? (\textbf{Seven}.) Which other character starts with the same piece and is said zuò? (\textbf{\zh{做}}.) What are the last two strokes? (\textbf{Two short bars on the right of the vertical}.)

Source: \href{https://raw.githubusercontent.com/chanind/hanzi-writer-data/68d10a4b21150cae5e1ebbd223eed289cf32d90c/data/\%E4\%BD\%9C.json}{Hanzi Writer Data \zh{作}, pinned stroke paths and medians}.
\end{tcolorbox}
\section[\texorpdfstring{\zh{写作} (xiězuò)}{写作 (xiězuò)}]{\zh{写作} (xiězuò) — to write essays and stories; writing}

\label{lesson:ZH-C251-xiezuo}



\begin{quote}
Before the new one: say the Mandarin for measure word for essays, then the Mandarin for a school essay.
\end{quote}

\subsection*{You'll want to know: \zh{写作}}
\textbf{\zh{写作}} — \emph{xiězuò} — \textquotedblleft{}to write essays and stories; writing\textquotedblright{}.

\begin{quote}
\textbf{\zh{我很喜欢写作}} \emph{wǒ hěn xǐhuan xiězuò} — I really enjoy writing.
\end{quote}

The new character, in pieces: \zh{写} — five strokes: a cover of a dot and a bar that turns down; under it a short bar, a stroke that comes down, turns right, turns down again and hooks, and a long bar at the foot. The next lesson puts it in your hand.

\begin{cousinweb}
\textbf{\zh{写作}} (\emph{xiězuò}) is built from \zh{写} 'to write' and \zh{作} 'to make', both inherited from Middle Chinese, \emph{sjaeX} and \emph{tsak} (Guangyun, in Baxter's transcription), which glosses the first, among other senses, 'to copy out a book'.
\end{cousinweb}

\begin{grammarlens}[title={What you've built}]
One more word for saying what I do.
\end{grammarlens}

\subsection*{Your turn}
\begin{itemize}
\raggedright
  \item \emph{Say it:} \emph{xiězuò}
  \item \emph{Say it:} \emph{xiězuò}, once more
  \item \emph{Say it:} \emph{wǒ hěn xǐhuan xiězuò} — I really enjoy writing.
  \item \emph{Recall:} say the Mandarin for measure word for essays, then the Mandarin for a school essay, then say \emph{xiězuò} again
\end{itemize}

\begin{tcolorbox}[breakable,colback=teal!4,colframe=teal!35!black,title={Before you move on}]
What does \zh{写作} mean? (\textquotedblleft{}To write essays and stories; writing\textquotedblright{}.) Say it once more.
\end{tcolorbox}
\section[\texorpdfstring{\zh{写} (xiě)}{写 (xiě)}]{\zh{写} — a cover, then a turning hook}

\label{lesson:ZH-W251-xie-write}



\begin{quote}
The word first, then the character inside it.

\begin{itemize}
\raggedright
  \item \emph{Recall:} say \emph{xiězuò} — \textbf{R1}, one lesson back, before it decays
\end{itemize}

\begin{quote}
\zh{学} \zh{字}
\end{quote}

\emph{Read it:} each one aloud, with nothing said first — \emph{xué}, \emph{zì}
\end{quote}

\subsection*{Script — a cover, then a turning hook}
\hlblockfigure{\detokenize{figures/ZH-W251-xie-write-filmstrip.pdf}}{How \zh{写} is written, stroke by stroke}

\begin{quote}
\zh{写}
\end{quote}

\textbf{Two pieces, top before bottom.}

\textbf{The cover, two strokes.} A short dot on the left, then a bar that runs right and turns down, ending in a small hook: the cover of \zh{学}.

\textbf{Under the cover, three strokes.} A short bar in the middle. Then a stroke that starts right under the cover, comes down, turns right, turns down again and hooks to the left. Then the long bar across the foot.

2 + 3 = five.

\subsection*{Your turn}
\begin{itemize}
\raggedright
  \item \emph{Write it:} \zh{写}, all five, with the model in front of you
  \item \emph{Cover:} write it again from memory, and count the pieces, not the strokes
\end{itemize}

\begin{tcolorbox}[breakable,colback=teal!4,colframe=teal!35!black,title={Before you move on}]
How many strokes in \zh{写}? (\textbf{Five}.) Where have you made the cover before? (\textbf{On \zh{学}, under its dots}.) What comes last? (\textbf{The long bar at the foot}.)

Source: \href{https://raw.githubusercontent.com/chanind/hanzi-writer-data/68d10a4b21150cae5e1ebbd223eed289cf32d90c/data/\%E5\%86\%99.json}{Hanzi Writer Data \zh{写}, pinned stroke paths and medians}.
\end{tcolorbox}
\section[\texorpdfstring{\zh{生字} (shēngzì)}{生字 (shēngzì)}]{\zh{生字} (shēngzì) — new character; a character not yet learned}

\label{lesson:ZH-C251-shengzi}



\begin{quote}
Before the new one: say the Mandarin for a school essay, then the Mandarin for to write essays and stories.
\end{quote}

\subsection*{You'll want to know: \zh{生字}}
\textbf{\zh{生字}} — \emph{shēngzì} — \textquotedblleft{}new character; a character not yet learned\textquotedblright{}.

\begin{quote}
\textbf{\zh{这些生字我都会写}} \emph{zhèxiē shēngzì wǒ dōu huì xiě} — I can write all these new characters.
\end{quote}

\begin{cousinweb}
\textbf{\zh{生字}} (\emph{shēngzì}) is built from \zh{生} and \zh{字} 'a character', both inherited from Middle Chinese, \emph{srjaeng} and \emph{dziH} (Guangyun, in Baxter's transcription), which glosses the first 'to grow'. Its sense 'raw, not yet known' is not given there, so how it came is not settled.
\end{cousinweb}

\begin{grammarlens}[title={What you've built}]
One more word for saying what I do.
\end{grammarlens}

\subsection*{Your turn}
\begin{itemize}
\raggedright
  \item \emph{Say it:} \emph{shēngzì}
  \item \emph{Say it:} \emph{shēngzì}, once more
  \item \emph{Say it:} \emph{zhèxiē shēngzì wǒ dōu huì xiě} — I can write all these new characters.
  \item \emph{Recall:} say the Mandarin for a school essay, then the Mandarin for to write essays and stories, then say \emph{shēngzì} again
\end{itemize}

\begin{tcolorbox}[breakable,colback=teal!4,colframe=teal!35!black,title={Before you move on}]
What does \zh{生字} mean? (\textquotedblleft{}New character; a character not yet learned\textquotedblright{}.) Say it once more.
\end{tcolorbox}
\section[xiūgǎi]{xiūgǎi — to revise; to correct (a piece of work)}

\label{lesson:ZH-C251-xiugai}



\begin{quote}
Before the new one: say the Mandarin for to write essays and stories, then the Mandarin for a new character.
\end{quote}

\subsection*{You'll want to know: xiūgǎi}
\emph{xiūgǎi} — \textquotedblleft{}to revise; to correct (a piece of work)\textquotedblright{}.

\begin{quote}
\emph{zhè piān zuòwén nǐ yào xiūgǎi} — You'll need to revise this essay.
\end{quote}

Its characters come later; for now it is a word for the ear.

\begin{cousinweb}
\emph{xiūgǎi} is built from \emph{xiū} and \emph{gǎi}, both inherited from Middle Chinese, \emph{sjuw} and \emph{kojX} (Guangyun, in Baxter's transcription), which glosses the first 'to put in order' and 'to adorn', and the second 'to change'.
\end{cousinweb}

\begin{grammarlens}[title={What you've built}]
One more word for saying what I do.
\end{grammarlens}

\subsection*{Your turn}
\begin{itemize}
\raggedright
  \item \emph{Say it:} \emph{xiūgǎi}
  \item \emph{Say it:} \emph{xiūgǎi}, once more
  \item \emph{Say it:} \emph{zhè piān zuòwén nǐ yào xiūgǎi} — You'll need to revise this essay.
  \item \emph{Recall:} say the Mandarin for to write essays and stories, then the Mandarin for a new character, then say \emph{xiūgǎi} again
\end{itemize}

\begin{tcolorbox}[breakable,colback=teal!4,colframe=teal!35!black,title={Before you move on}]
What does xiūgǎi mean? (\textquotedblleft{}To revise; to correct (a piece of work)\textquotedblright{}.) Say it once more.
\end{tcolorbox}
